\section{Постановка задачи}
\label{sec:task}

Целью лабораторной работы является изучение принципов работы и проектирования решателей задач интеллектуальных систем, освоение навыков проектирования и реализации программ с помощью различных видов представления знаний и моделей решения задач, научиться анализировать и выводить знания.

Задача работы состоит в разработке модуля адаптации знаний для системы NIKA: агента, который получает знания из внешнего источника, погружает их в базу знаний и позволяет системе отвечать на вопросы пользователя даже при отсутствии описывающих фактов в локальной БЗ.

Вариант~\mdash  15. Предметная область~\mdash  музыка (класс сообщений и сущности заданы в лабораторной работе №2, логические правила и фразы~\mdash  в лабораторной работе №3). Класс сообщения \texttt{concept\_message\_about\_entity} относится к сообщениям вида <<Что такое ...?>>, <<Кто такой ...?>>, <<Расскажи о ...>>.

Индивидуальный вариант~\mdash  вектор экземпляров категорий из четырёх классов (таблица~\ref{tab:variant}).

\begin{table}[H]
	\centering
	\small
	\begin{tabular}{|p{4.6cm}|p{4.6cm}|p{5.2cm}|}
		\hline
		Класс & Выбранное значение & Обоснование \\
		\hline
		Локальность источника & Внешние сервисы, доступные по сети Internet & Агент обращается к MediaWiki API Википедии по протоколу HTTPS \\
		\hline
		Структура внешней информации & Полуструктурированные данные: JSON & Ответы MediaWiki API представляются в формате JSON \\
		\hline
		Модель обработки знаний & Процедурные (алгоритмические) модели & Обработка~\mdash  пошаговый алгоритм: получить сущность, имя, запрос, разобрать ответ, записать факты \\
		\hline
		Язык программирования & Python & Агент реализован как Python-модуль \texttt{sc-kpm} \\
		\hline
	\end{tabular}
	\caption{Индивидуальный вариант лабораторной работы}
	\label{tab:variant}
\end{table}

Требуется:
\begin{enumerate}
	\item спроектировать и реализовать агента, получающего знания из внешнего источника и погружающего их в БЗ;
	\item обеспечить вывод новых знаний на основе погруженных: на добавленной заметке должно строиться описание сущности в ответе системы;
	\item интегрировать агента с общей программой обработки сообщений NIKA;
	\item протестировать работу агента средствами диалога с системой;
	\item задокументировать агента: условие инициирования, результат, входы, выходы, процедуры и схему алгоритма.
\end{enumerate}
